\scanpage{12}

\mspage{6}

\uline{Théorème 1} Supposons que $M$ soit \uline{accessible}, et \deleted{que les sommes \unsure{amalgamées|quelconques|inductives}, et par} \inserted{stable par (petites) limites inductives}. (Il doit suffire stabilité par sommes amalgamées, et par $\varinjlim$ filtrantes, voire ordinales\ldots)

\deleted{\unsure{a)}} Soit $\Phi_0$ une \uline{petite} partie de $\mathrm{Fl}\,M$, \deleted{et $\Phi$ une} $\Psi=(\Phi_0)_*$, et $\Phi$ une partie de $\mathrm{Fl}\,M$ telle que

\[ \Phi_0\subset\Phi\subset\Psi^*\ (= \widetilde{\Phi}_0). \]

Alors a) Pour que $\Phi$ soit \deleted{\unsure{stable}} \inserted{close} à gauche, \inserted{i.e. $\Phi=\widetilde{\Phi}_0$}, il f. et s. qu'\unsure{elle|elles} \unsure{satisfasse|satisfassent} les conditions a) b),c), de la prop.~2. Quand il en est ainsi, le couple $(\Phi,\Psi)$ est un couple de Quillen.

b) Supposons seulement que $\Phi$ satisfasse aux conditions a), c) de la prop 2 (à l'exclusion de la condition b) des facteurs directs) -- alors on ne peut affirmer que $\Phi$ soit close à gauche. Néanmoins, $(\Phi,\underset{\Phi_*}{\underset{\Vert}{\Psi}})$ satisfait à la propriété de factorisation de Quillen.

\uline{Remarque} Dans la démonstration, l'existence de \deleted{\unsure{un}} $\Phi_0$\inserted{$\subset\Phi$,} qui joue le rôle d'un petit ens. de ``générateurs\unsure{''} \deleted{pour} (à g.) pour $\Phi$, joue un rôle essentiel. Il y aurait lieu de développer la notion \unsure{des|de} \uline{parties accessibles}

% archive sheet number: 11

